% ===================================================================
%  Shared TikZ styles for the schematic figures (Figs. 1-3).
%
%  The palette is the same semantic palette as figs/style.py, so the
%  schematics and the measured panels read as one figure set:
%
%     fgBlue  #1F5FA9   the learned / adopted path      (style.py "learned")
%     fgRed   #C1272D   the oracle and the failed claim (style.py "oracle")
%     fgGrey  #4A4A4A   the control path, neutral       (style.py "periodic")
%
%  Everything is drawn at true size in millimetres (x=1mm,y=1mm) and each
%  picture is <= 84 mm wide, which is \columnwidth for acmart sigconf at
%  10 pt. Nothing is ever passed through \resizebox: scaling is what made
%  the previous versions of these figures illegible.
%
%  Smallest type used anywhere below is \scriptsize = 7 pt.
% ===================================================================

\definecolor{fgInk}{HTML}{2B2B2B}
\definecolor{fgMuted}{HTML}{6E6E6E}
\definecolor{fgBlue}{HTML}{1F5FA9}
\definecolor{fgRed}{HTML}{C1272D}
\definecolor{fgGrey}{HTML}{4A4A4A}

\tikzset{
  % ---------------------------------------------------------- containers
  band/.style      = {rounded corners=1.1pt, draw=none},
  bandlbl/.style   = {font=\scriptsize\bfseries, inner sep=0pt},
  % ---------------------------------------------------------- boxes
  bx/.style        = {rounded corners=1.1pt, draw=fgInk!45, line width=0.4pt,
                      fill=white, align=center, font=\scriptsize,
                      inner sep=0.7mm, text=fgInk},
  bxb/.style       = {bx, draw=fgBlue!70, fill=fgBlue!7},
  bxr/.style       = {bx, draw=fgRed!70,  fill=fgRed!7},
  bxg/.style       = {bx, draw=fgGrey!60, fill=black!5},
  % ---------------------------------------------------------- flow arrows
  ar/.style        = {-{Latex[length=1.35mm,width=1.05mm]},
                      draw=fgInk!70, line width=0.45pt},
  arb/.style       = {ar, draw=fgBlue},
  arr/.style       = {ar, draw=fgRed},
  arg/.style       = {ar, draw=fgGrey!75},
  arfb/.style      = {ar, draw=fgGrey!70, dashed, dash pattern=on 0.9mm off 0.7mm},
  % ---------------------------------------------------------- text
  lbl/.style       = {font=\scriptsize, text=fgMuted, inner sep=0.35mm,
                      align=center},
  lblink/.style    = {lbl, text=fgInk},
  lblb/.style      = {lbl, text=fgBlue},
  lblr/.style      = {lbl, text=fgRed},
  tag/.style       = {font=\scriptsize\bfseries, inner sep=0.35mm,
                      align=center},
  rule/.style      = {draw=fgInk!22, line width=0.35pt},
}

% ------------------------------------------------------------------ glyphs
% A minimal camera, drawn rather than imported: the previous figures used
% clip art whose line weight and grey did not match anything else on the page.
% #1 = centre, #2 = draw colour.
\newcommand{\camglyph}[2]{%
  \colorlet{glyphc}{#2}%
  \begin{scope}[shift={#1}]
    \fill[glyphc!12] (-3.1,-1.8) rectangle (2.1,1.8);
    \draw[draw=glyphc, line width=0.4pt] (-3.1,-1.8) rectangle (2.1,1.8);
    \draw[draw=glyphc, line width=0.4pt, fill=white] (-1.0,0) circle (0.95);
    \draw[draw=glyphc, line width=0.4pt, fill=glyphc!25]
        (2.1,1.0) -- (3.5,1.8) -- (3.5,-1.8) -- (2.1,-1.0) -- cycle;
  \end{scope}}

% A padlock, used in Fig. 3 to mark each irreversible freeze.
% #1 = centre, #2 = draw colour.
\newcommand{\lockglyph}[2]{%
  \colorlet{glyphc}{#2}%
  \begin{scope}[shift={#1}]
    \draw[draw=glyphc, line width=0.4pt] (-0.62,0.28) arc (180:0:0.62);
    \draw[draw=glyphc, line width=0.35pt, fill=glyphc!16, rounded corners=0.15pt]
        (-1.02,-1.02) rectangle (1.02,0.32);
  \end{scope}}

% Verdict badges for Fig. 2. Drawn rather than set from a symbol font so the
% stroke weight matches the rules and arrows around them.
% #1 = centre, #2 = colour.
\newcommand{\okglyph}[2]{%
  \colorlet{glyphc}{#2}%
  \begin{scope}[shift={#1}]
    \draw[draw=glyphc, line width=0.4pt, fill=glyphc!10] (0,0) circle (1.35);
    \draw[draw=glyphc, line width=0.55pt, line cap=round, line join=round]
        (-0.62,0.05) -- (-0.18,-0.5) -- (0.68,0.6);
  \end{scope}}
\newcommand{\noglyph}[2]{%
  \colorlet{glyphc}{#2}%
  \begin{scope}[shift={#1}]
    \draw[draw=glyphc, line width=0.4pt, fill=glyphc!10] (0,0) circle (1.35);
    \draw[draw=glyphc, line width=0.55pt, line cap=round]
        (-0.52,-0.52) -- (0.52,0.52)  (-0.52,0.52) -- (0.52,-0.52);
  \end{scope}}
